\scanpage{31}\mspage{30}
% @ top, part a
a) $F'$ est une figure. \deleted{\unsure{i.e.}} Notons que \deleted{$|F'|=$}
% @ top, support with braces
\[
 |F'|=\underbrace{|G'|}_{|G|}\cup\underbrace{|F\setminus G|}_{\supset |F|\setminus |G|}\deleted{\unsure{\supset|=}}\supset |F|
\]
% @ upper part, connector
donc, comme $|F'|\subset |F|$, on a
% @ upper part, support equality
\[
 |F'|\unsure{=|\supseteq}|F|
\]
% @ upper part, argument
Il faut prouver que $\forall x\in |F'|=|F|$, $\exists$ plus petit \unsure{élémt|élém\textsuperscript{t}|élément|élt} de $F'$ qui le contienne. Distinguons deux cas

% @ upper middle, first case
1$^\circ$) $x\in |G|=|G'|$, soit $A\in G'\deleted{\unsure{\circ}}$ tel que $x\in A^{\circ(G')\deleted{\unsure{)|?}}}$. Si $B\in F'$ et \deleted{$\unsure{B\supset|B}$} $x\in B$, prouvons $A\subset B$. Si $B\in G'$, c'est évident. Si $B\in F\setminus G$, alors \deleted{$B$ \unsure{contient}} $A^{\circ(G')}\cap B\ne\phi$ donc $A^{\unsure{\circ|\bullet|*}}\subset B$ OK

% @ middle, second case
2$^\circ$) $x\in |\deleted{\unsure{G}}F|\setminus |G|$, alors les $\deleted{\unsure{B}}A\in F'$ qui contiennent $x$ ne sont pas dans $G'$, ils sont donc dans $F\setminus G$. D'autre part, l'unique \unsure{élémt|élém\textsuperscript{t}|élément|élt} $A$ de $F$ tel que $x\in A^{\circ(F)}$ n'est pas dans $G$ puisque $x\notin |G|$, il est donc dans $F\setminus G$, et \deleted{\unsure{si $x\in A$}} c'est lui plus petit él. de $F$ et à fortiori de $F\setminus G$) qui contient $x$, OK~!

% @ lower middle, part b
b) $F'$ \deleted{\unsure{raffine}}\inserted{subdivise} $F$. Comme $|F'|=|F|$, il faut voir \unsure{seult|seulement} que $F'$ raffine $F$. Donc \deleted{que \unsure{tout}} que
% @ lower part, cancelled attempt
\deleted{1$^\circ$) \deleted{$\forall$} $A\in F'$, \unsure{et} \inserted{l'ens $F^A$ des $B\in F$ tels que $A\subset B$} \unsure{est contenu dans|contienne} un plus petit}
% @ one continuous source cancellation; wrappers split only to place the cancelled insertion in the margin
\margin{\deleted{\inserted{\unsure{a un plus petit|$=$ un plus petit}}}}
\deleted{\unsure{élémt|élém\textsuperscript{t}|élément|élt} de $F$. C'est \unsure{évident} si $A\in F$ c'est \unsure{évident}. Si}
\deleted{\unsure{str.} de de $F'$ \unsure{raffine} $F$. Si}

% @ lower part, restated first condition
1$^\circ$) \deleted{$\forall$} $A\in F'$, l'ens $F^A=\{B\in F\mid \deleted{\unsure{B}}A\subset B\}$ a un plus petit \unsure{élémt|élém|élément|élt}. C'est évident si $A\in F$ \deleted{$(B=A)$} (on prend $B=A$) et si $A\in G'$, c'est vrai puisque $G'\ll F$.

% @ bottom, second condition
2$^\circ$) Si $A\in F'$, $B\in F$, alors $A\cap B$ est \deleted{\unsure{$\mathcal R_{F'}$-satu.|$\mathscr R_{F'}$-satu.|$Q_{F'}$-satu.|$\mathcal O_{F'}$-satu.}}\inserted{fermé dans $F'$}, i.e. réunion \unsure{de|des} strates de $F'$. Si $A\in G'$, comme $G'\ll F$, $A\cap B$ est réunion de strates de $G'$, qui sont aussi des strates de $F'$.
